\documentclass[11pt]{article}
\usepackage[margin=1in]{geometry}
\usepackage[T1]{fontenc}
\usepackage{lmodern,amsmath,amssymb,amsthm,booktabs,microtype}
\usepackage[hidelinks]{hyperref}
\newtheorem{theorem}{Theorem}
\setlength{\emergencystretch}{3em}
\allowdisplaybreaks
\title{Rectangular common-basis geometry for integer multiplication\\\large Conditional saving $12260937/10^{12}$}
\author{Rohan Arun\\\small With substantial OpenAI Codex assistance}
\date{8 October 2026}
\begin{document}\maketitle
\begin{abstract}
We extend the retained partial-swap common-basis construction to tensor dimensions $(28,27,57)$ and exhibit a contiguous width-25 block inside its A5 corner. The resulting exact bit saving is $12261239/10^{12}$. Composing with the inherited semantic/bulk interfaces gives the conditional bound
\[
T(n)=O\bigl(n(\log n)^{1-\kappa}\bigr),\qquad
\kappa=\frac{12260937}{10^{12}}=1.2260937\times10^{-5}.
\]
The changed bit network requires joint row degree $96000$. Exact rational certificates and finite controls support, but do not formally verify, the written geometric argument and inherited analytic/tape hypotheses.
\end{abstract}

\section{Scope and retained construction}
The physical positive producers and prescribed partial-swap bases are due to icekylinx's PR18; translated source frames are Zhihao Chen's PR21. We use Zhihao Chen's PR23 composition of RaD/PR20 semantic precision and bulk resampling. The earlier A5 observation in PR25 is generalized here. We credit eumemic's source-frame contribution and all earlier authors listed in \texttt{NOTICE}. The new contribution is rectangular compatibility, the conservative width-25 block, fresh dimension-specific producer evidence and their exact assembly. This is not a claim of practical speedup, global optimality, or removal of any inherited hypothesis.

Let $a=28,b=27,k=57$, $H=ab=756$, $m=Hk=43092$, $s=k-2a=1$, and $n=a+k=85$. With $I=(0,\ldots,a-1)$ and $S=(0,\ldots,s-1)$, retain the prescribed row labels $R=[I,S,I,I]$ and inverse-column labels $C=[I,I,S,I]$. At each residue modulo $b$, the prescribed entries extend to a permutation. The exact certificate lists those completions. Both bipartite nullspace graphs on $a+b$ vertices are trees, as verified by their explicit edges and connectivity. Also $2n\leq H$, $H>2k$, and both $a,k$ are nonzero modulo $b$.

Fix these permutation completions. The remaining invertible inner bases and outer completion matrices form irreducible rational open families. A finite collection of individually nonzero minor polynomials therefore admits one common rational specialization. This step requires nonvanishing in this prescribed family, not merely in a family of arbitrary dense matrices. We establish the required A1, A3 and A5 conditions below. The ordinary local profiles and translated exit retain their matched first/last prescriptions and the corresponding inherited nonzero-minor arguments.

\section{A1: structural separation of coefficient labels}
Write normalized line pairs as $p\xi$, $t\phi$, and two orthogonal inner pairs $u\eta,v\nu$, so $\eta u=\nu v=1$, $\eta v=\nu u=0$. Put
\[
o=[0,\ldots,k-1,0,\ldots,a-1],\quad
o'=[0,\ldots,a-1,0,\ldots,k-1],
\]
$\beta_i=i\bmod b$ and $\beta'_j=(H-n+j)\bmod b$. The actual factorized null-projector corner is
\[
Q_{ij}=\mathbf1_{o_i=o'_j}p_{R_i}\xi_{C_j}v_{\beta_i}\nu_{\beta'_j}
 +\mathbf1_{R_i=C_j}t_{o_i}\phi_{o'_j}u_{\beta_i}\eta_{\beta'_j}.
\]
Equality $o_i=o'_j$ implies $R_i=C_j$, so elimination preserves the direct sum by coefficient label. Each label below $s$ contributes a $4\times4$ restriction; each other label contributes a $3\times3$ restriction. Within each restriction the first summand has rank one on a duplicated outer coordinate and nonzero scalar blocks on the others. The second summand is rank one. Nonsingularity requires its row and column vectors to escape the respective spaces of the first summand. These conditions compare $u_\beta/v_\beta$ at coordinates shifted by $k\bmod b$ and $\eta_\beta/\nu_\beta$ at coordinates shifted by $a\bmod b$. These are distinct coordinates; the dual-line orthogonality equations do not force the ratios equal. Their nonvanishing conditions are nonzero rational functions.

Ordered rightmost-pivot elimination first uses the last $a$ columns against the first $a$ rows, a diagonal block from the second summand. The next $s$ pivots are diagonal scalar entries of the first summand after cancelling the rank-one term. The remaining corner is a $2\times2$ matrix of diagonal $a$-blocks. Its upper-right entries are nonzero rational functions: after denominator clearing, setting the relevant outer-line coordinates to zero leaves a product of nonzero first-summand scalar entries; normalize the outer line at an unused coordinate. For a $4\times4$ restriction also set its extra outer-row coordinate to zero, reducing to the same test while retaining its scalar pivot. Such specializations establish nonzero polynomials and need not satisfy the other open conditions simultaneously. The final pivots follow from full nonsingularity. The retained large-projector identity thus gives
\[
A1:\quad [a,s,a,a,m-2n]=[28,1,28,28,42922].
\]

\section{A3: the matching interval and the full nullity corner}
Let $P_q=q\xi$ be the inner rank-one projector in the retained basis, with nonzero coordinates, and $P_t=t\phi$ the outer projector. The A3 projector is $(I_k-P_t)\otimes(I_H-P_q)$, of nullity $r=H+k-1$. The controlled family acts as an invertible $G_i$ at multiplicity $i$. The top-$H$/rightmost-$H$ corner has the exact form
\[
D-\operatorname{diag}(q)UV\operatorname{diag}(\xi),\quad
D_i=-(G_it)_0(\phi G_i^{-1})_{k-1},\quad
U_i=e_0^TG_i(I-P_t),\quad V_j=G_j^{-1}e_{k-1}.
\]
The rank perturbation has rank $k-1$. Matched first-$k$ rows and last-$k$ inverse columns give a diagonal rescaling of $I-L_0P_tL_0^{-1}$ on their disjoint crossing. The ordinary rank-$(k-1)$ lemma supplies one singleton and a $(k-2)$-block, using rows $0,\ldots,k-2$ and columns $H-k+1,\ldots,H-1$. They exhaust the perturbation. On the untouched matching interval $[k-1,H-k]$, all crossings to these selected coordinates have zero contribution from $D$. Its Schur complement is therefore precisely $D$ on that interval, supplying $H-2k+2$ consecutive diagonal pivots. Remaining boundary columns lie to their left.

Each $D_i$ is generically nonzero in the prescribed family. At the first boundary only the first row is fixed; at the last boundary only the last inverse column is fixed. The two boundaries are disjoint. A free inverse column must annihilate the prescribed first row, but may pair nontrivially with $\phi$ unless that row is proportional to $\phi$. The dual statement requires $t$ not proportional to the prescribed inverse column. These proper conditions are avoided by the free $L_0$ together with nonzero prescribed line pairings.

It remains essential to prove full nullity-corner rank. A nullspace vector before control has the form $w\otimes q+t\otimes z$, modulo the one-dimensional redundancy $(w,z)=(ct,-cq)$. Its controlled coordinates are
\[
q_i(G_iw)_\gamma+(G_it)_\gamma z_i.
\]
The first $H$ rows determine $z_i=-q_i(G_iw)_0/(G_it)_0$. In the next $k-1$ rows ($\gamma=1$, $i=0,\ldots,k-2$), the residual functionals of $w$ are
\[
q_i\left[e_1^TG_i-\frac{(G_it)_1}{(G_it)_0}e_0^TG_i\right].
\]
Only the first row is prescribed there. Choose the free second rows to form a basis of $\operatorname{Ann}(t)$; each is independent of its prescribed first row, which has nonzero value on $t$. Complete each matrix invertibly. The displayed functionals then force $w\parallel t$, establishing a nonzero left-restriction minor. Dually, the penultimate inverse columns at the opposite boundary are free and yield a nonzero right-restriction minor. These minors, the matching pivots and all previous conditions coexist by irreducibility of the same completion family.

For a null-projector factorization $Q=XY$ with $YX=I_r$, the full disjoint boundary corner is $-X_{\rm left}Y_{\rm right}$ and hence invertible. Retain the other $2k-2$ corner pivots as singletons and apply the large-projector identity to the middle. Altogether,
\[
A3:\quad 113\text{ singletons}+ [55,644,41468].
\]

\section{A5: a contiguous width-25 block}
The A5 nullity is $r_5=a+b-1=54$. Its row-label prefix is $[0..27,0,0..24]$ and column-label suffix is $[3..27,0,0..27]$. As $H-r_5=702$ is divisible by $27$, its actual null corner is
\[
Q_{ij}=\mathbf1_{i\bmod27=j\bmod27}p_{R_i}\xi_{C_j}
 +\mathbf1_{R_i=C_j}v_{i\bmod27}\nu_{j\bmod27}
 -p_{R_i}\xi_{C_j}v_{i\bmod27}\nu_{j\bmod27}.
\]
The first ordered pivot $(0,53)$ is the nonzero pure rank-one entry $-p_0\xi_{27}v_0\nu_{26}$. Its Schur update cancels the final term on rows $i=1,\ldots,25$, columns $j=28,\ldots,52$. Write $j=27+\ell$. The remaining entries are exactly
\[
\mathbf1_{i=\ell}p_i\xi_{\ell+1}
 +\mathbf1_{i=\ell+1}v_i\nu_\ell.
\]
This is lower bidiagonal with nonzero diagonal. Thus each successive row $i$ pivots at column $27+i$, all available entries to its right being zero. Earlier eliminations preserve these pivots. They are contiguous in physical order and require no gathering.

After diagonal rescaling, both full nullspace restrictions evaluate $x_{R_i}+y_\beta$. Their prescribed graphs are trees, so leaf elimination proves rank $a+b-1$ modulo the constant redundancy. All remaining 29 corner pivots may conservatively stay singleton. The separate middle block has width $H-2r_5=648$, giving
\[
A5:\quad29\text{ singletons}+[25,648].
\]
The null corner sign for the actual projector does not change the profile. These identities and tree conditions hold in the same prescribed family as A1/A3, so their nonzero minor conditions can be multiplied together.

\section{Counts, precision and exact assembly}
The unchanged producer algorithms are freshly replayed at $h=27,28$; the retained $h=57$ has a fresh PR25 full-run receipt and is regenerated by the inherited full verification target. Their positive role counts are respectively $66258,75292,729786$. Original dependency matching, positive labels, positive matching and exact rank histograms are checked separately. The dimension-parametric physical construction gives
\[
\begin{split}
N&=280378098000,\quad L=162580083120,\\
(B_1,B_2,B_3)&=(6443903466000,6351210946080,6993028387800),\\
W&=13904995529880,\quad m=43092,\quad M=43038,\\
s_b&=599193831777559200= Wm-2N+2L.
\end{split}
\]
There are $B_1$ A1 occurrences, $2N$ each of A3 and A5, and $B_2$ translated exits with profile $[27,43038]$. The remaining ordinary/growth profiles and every width multiplicity are explicitly reconstructed in the certificate. Replacing 25 singletons by one A5 child preserves rank mass. The exact log/exponential upper enclosure gives
\[
\frac1W\sum_t c_t(t/m)^{1-a_b}<1,\qquad a_b=12261239/10^{12}.
\]
Logs use rational series with positive tail bounds; the exponential uses $e^u\leq1+u+u^2/(2(1-u/3))$ for $0\leq u<1$. The next grid point fails this particular sufficient enclosure, which does not establish impossibility or optimality.

The complex interface is unchanged, with saving $a_c=18/10^6$, dimension $21952$, and maximum child $21924$. Its semantic constants, literal charge, guard $C_1=1$, and scalar induction are identical to pinned PR23/25. The new bit wire count needs 44 bits and halving degree553; the complex values remain41 and544. The joint row coefficient is $44\cdot553+41\cdot544=46636$. Therefore use degree96000 and suffix slope384000; the strict degree margin against $(51/25)46636$ is $21564/25$. This increased stock is included in every cutoff and constraint.

Put $h=10^{-8}$, $\beta=1/4$, $\tau=1-a_b$, $\sigma=1-a_c$,
\[
q=a_b(1-2h),\quad c=q(1+h),\quad
\epsilon=\frac{1-h}{1+c+q},\quad
\lambda'=1-q,\quad \lambda=\frac{\tau+\lambda'}2,
\]
\[
G=\epsilon q,\quad r=\frac{G+1-\epsilon}{2},\quad\delta=h/8.
\]
The standalone exact checker recomputes all47 strict inequalities and seven margins, including their explicit eventual cutoffs. Their minimum is $G$ and $G-\kappa>9\cdot10^{-13}$. Negative controls reject the old quadratic guard, old separate exposure bounds, and the next $10^{-12}$ grid point for $\kappa$. No inherited analytic, Gaussian inverse, uniformity, ordered-affine or fixed finite-alphabet multitape hypothesis is removed.

\section{Reproduction and review boundary}
Run \texttt{make rectangular-semantic-certificate rectangular-semantic-patch} for the exact record and focused source patch; \texttt{make rectangular-semantic-producers} regenerates the new finite producers. Full \texttt{make verify} additionally reruns inherited checks. The validation receipt distinguishes targeted checks from completion of the full rerun. The focused patch applies to pinned PR25 commit \texttt{003f366}, not to the original OpenAI manuscript. Inherited integrated patches remain unchanged. Exact controls test actual A1/A5 corners, arbitrary rational line choices and deliberate wrong-boundary/omitted-cleanup variants; they supplement the general proof rather than replace it.

\begin{thebibliography}{9}
\bibitem{pr18} icekylinx, \href{https://github.com/CrocSwap/integer-mult-bounds/pull/18}{PR18}, partial-swap construction and prescribed basis.
\bibitem{pr21} Zhihao Chen, \href{https://github.com/CrocSwap/integer-mult-bounds/pull/21}{PR21}, translated source frames.
\bibitem{pr20} RaD, \href{https://github.com/CrocSwap/integer-mult-bounds/pull/20}{PR20}, semantic/bulk interfaces.
\bibitem{pr23} Zhihao Chen, \href{https://github.com/CrocSwap/integer-mult-bounds/pull/23}{PR23}, semantic composition.
\bibitem{pr25} Rohan Arun with OpenAI Codex assistance, \href{https://github.com/CrocSwap/integer-mult-bounds/pull/25}{PR25}, retained A5 block.
\end{thebibliography}
\end{document}
